\section{Experiments}
\label{sec:experiments}

\paragraph{Datasets.}
We evaluate on three multi-batch single-cell RNA-seq datasets~\citep{luecken2022benchmarking}:
\textbf{Pancreas} (16,382 cells, 9 sequencing technologies used as batches, 14 cell types),
\textbf{PBMC (immune)} (30k cells, 5 batches, 9 cell types),
and a \textbf{Lung} dataset (30k cells, 16 batches, 17 cell types).
For spatial evaluations, we use a \textbf{non-small cell lung cancer (NSCLC)} spatial transcriptomics dataset~\citep{pentimalli2025combining}.
After preprocessing (see Appendix~\ref{app:datasets}), this yields 28,804 cells across 18 cell types and 9 niches.
Niche labels are defined by clustering cells based on the local cell-type composition of their spatial neighborhood, providing a measure of microenvironmental context that is related to but distinct from the spatial affinity used during training.

\paragraph{Affinity construction.}
For transcriptomic experiments, cell--cell affinities are computed using an adaptive Gaussian kernel on PCA embeddings~\citep{seacells}, which scales each cell's bandwidth to its local neighborhood density, ensuring rare cell types form well-connected communities despite occupying sparse regions; combined with scProto's cross-batch aggregation, this enables stable metacells for rare populations underrepresented in individual batches.
For spatial experiments, we build a BANKSY-style affinity~\citep{banksy}: each cell's own PCA embedding is concatenated with a distance-weighted mean of its spatial neighbors' PCA embeddings (its spatial context), balanced by a spatial weight $\lambda$, and the same adaptive Gaussian kernel is then applied to the concatenated features (see Appendix~\ref{app:spatial_affinity}).

\paragraph{Baselines.}
We compare against three \textbf{metacell} methods --- \textbf{SEACells}~\citep{seacells} (archetypal analysis), \textbf{MetaCell2}~\citep{metacell2} (divide-and-conquer partitioning), \textbf{SuperCell}~\citep{supercell} (graph coarse-graining) --- \textbf{MetaQ}~\citep{metaq} (end-to-end vector-quantized, no batch correction), \textbf{scPoli (Stage-1)} (scProto's pretrained encoder before prototypes, Section~\ref{sec:method}), and three \textbf{batch-correction} methods --- \textbf{Harmony}~\citep{harmony}, \textbf{BBKNN}~\citep{bbknn}, \textbf{ComBat}~\citep{combat}.
All methods use $d=8$; scPoli (Stage-1) and the batch-correction methods lack native clusters, so we run Leiden~\citep{traag2019leiden} on their embedding, tuned to match scProto's $K$ (the rest take $K$ directly). For spatial experiments we add \textbf{SEACells (spatial)}: its default PCA affinity replaced by the spatial-context affinity.

\paragraph{Evaluation protocol.}
scProto outputs soft assignments $s_i \in \mathbb{R}^K$; cells take the argmax prototype and metacells are labeled by majority vote. We report \textbf{cohesion} (within-metacell similarity, Appendix~\ref{app:cohesion}), \textbf{batch entropy}, \textbf{cell-type purity}, rare-cell \textbf{coverage}/\textbf{homogeneity}, and, for spatial data, \textbf{niche purity}/\textbf{DGE} (formal definitions in Appendix~\ref{app:metrics}).
Table~\ref{tab:graph_batch} additionally uses a \textbf{node-holdout} protocol: 20\% of cells (stratified by batch) are excluded from the graph and training, with scProto/scPoli (Stage-1) placing them via a frozen-encoder forward pass, so those metrics are computed on never-trained-on cells only.
Rare types are identified within each batch (frequency more than one std below the batch mean, else lowest quartile); scores are mean~$\pm$~std across batches.

\subsection{scProto forms cohesive metacells while integrating cells across batches}

As shown in Table~\ref{tab:graph_batch}, scProto matches the metacell baselines' cohesion (essentially tied with SEACells, the strongest of them) while clearly exceeding every batch-effect-removal method on it; batch entropy shows the reverse pattern, with scProto comparable to the batch-effect-removal methods and clearly ahead of every metacell method, which mix nothing across batches by construction (e.g.\ SEACells: 0.16/0.15/0.38 vs.\ scProto's 1.21/0.70/1.16). Cell-type purity stays within 0.06 of every other method on Pancreas and Immune, so cross-batch alignment does not cost biological identity there; Lung is the exception, where scProto's purity (0.72) trails every metacell baseline by 0.18--0.19, a real dataset-specific cost. Fig.~\ref{fig:pancreas_umap} illustrates this: scProto's UMAP~\citep{mcinnes2018umap} clusters align across batches while staying consistent with cell type, whereas SEACells' are largely batch-separated.

\begin{figure*}[t]
\centering
\includegraphics[width=\textwidth]{figures/pancreas_umap.png}
\caption{
UMAP visualization of metacells on the Pancreas dataset (top: scProto, bottom: SEACells; left: cell type, right: batch). scProto forms compact, biologically coherent clusters while mixing cells across batches; SEACells produces batch-specific metacells. Circles with black outlines denote metacells (size $\propto$ assigned cells).
}
\label{fig:pancreas_umap}
\end{figure*}

\begin{table*}[t]
\centering
\caption{Node-holdout protocol (metric definitions in \S\ref{sec:experiments}): scProto/scPoli (Stage-1) place held-out cells via a frozen encoder; all values computed on held-out cells only. Other methods are fit on all cells (context, not matched). Bold: best; underlined: second-best.}
\label{tab:graph_batch}
\tiny
\setlength{\tabcolsep}{1pt}
\resizebox{\textwidth}{!}{%
\begin{tabular}{lcccc|cccc|cccc}
\toprule
 & \multicolumn{4}{c}{Pancreas} & \multicolumn{4}{c}{Immune} & \multicolumn{4}{c}{Lung} \\
\cmidrule(lr){2-5} \cmidrule(lr){6-9} \cmidrule(lr){10-13}
Method & Coh(aff) $\uparrow$ & Coh(gram) $\uparrow$ & Entropy $\uparrow$ & Purity $\uparrow$ & Coh(aff) $\uparrow$ & Coh(gram) $\uparrow$ & Entropy $\uparrow$ & Purity $\uparrow$ & Coh(aff) $\uparrow$ & Coh(gram) $\uparrow$ & Entropy $\uparrow$ & Purity $\uparrow$ \\
\midrule

scProto
& $\mathbf{0.18 \pm 0.13}$
& $\underline{0.33 \pm 0.19}$
& $\underline{1.21 \pm 0.40}$
& $0.93 \pm 0.12$
& $\underline{0.08 \pm 0.10}$
& $\underline{0.16 \pm 0.16}$
& $0.70 \pm 0.37$
& $0.82 \pm 0.16$
& $\underline{0.13 \pm 0.11}$
& $\underline{0.27 \pm 0.17}$
& $1.16 \pm 0.45$
& $0.72 \pm 0.19$ \\

scPoli (Stage-1) + Leiden
& $0.16 \pm 0.12$
& $0.28 \pm 0.16$
& $1.20 \pm 0.46$
& $\underline{0.97 \pm 0.08}$
& $0.06 \pm 0.06$
& $0.11 \pm 0.11$
& $\mathbf{0.88 \pm 0.43}$
& $\underline{0.87 \pm 0.14}$
& $0.10 \pm 0.10$
& $0.19 \pm 0.14$
& $\underline{1.41 \pm 0.59}$
& $0.84 \pm 0.20$ \\

\midrule

SEACells
& $\underline{0.18 \pm 0.11}$
& $\mathbf{0.36 \pm 0.18}$
& $0.16 \pm 0.28$
& $\mathbf{0.97 \pm 0.08}$
& $\mathbf{0.10 \pm 0.10}$
& $\mathbf{0.21 \pm 0.17}$
& $0.15 \pm 0.25$
& $\mathbf{0.88 \pm 0.15}$
& $\mathbf{0.14 \pm 0.10}$
& $\mathbf{0.30 \pm 0.16}$
& $0.38 \pm 0.45$
& $0.90 \pm 0.15$ \\

BBKNN
& $0.11 \pm 0.09$
& $0.21 \pm 0.14$
& $0.83 \pm 0.44$
& $0.92 \pm 0.12$
& $0.05 \pm 0.06$
& $0.10 \pm 0.10$
& $0.51 \pm 0.43$
& $0.83 \pm 0.15$
& $0.07 \pm 0.08$
& $0.15 \pm 0.12$
& $1.18 \pm 0.41$
& $0.80 \pm 0.20$ \\

Harmony
& $0.15 \pm 0.11$
& $0.25 \pm 0.15$
& $\mathbf{1.24 \pm 0.37}$
& $0.91 \pm 0.16$
& $0.08 \pm 0.06$
& $0.14 \pm 0.11$
& $\underline{0.84 \pm 0.45}$
& $0.84 \pm 0.15$
& $0.08 \pm 0.09$
& $0.15 \pm 0.13$
& $\mathbf{1.49 \pm 0.45}$
& $0.72 \pm 0.23$ \\

ComBat
& $0.17 \pm 0.12$
& $0.30 \pm 0.16$
& $0.82 \pm 0.48$
& $0.94 \pm 0.13$
& $0.06 \pm 0.08$
& $0.12 \pm 0.13$
& $0.42 \pm 0.43$
& $0.84 \pm 0.15$
& $0.10 \pm 0.11$
& $0.19 \pm 0.17$
& $1.00 \pm 0.59$
& $0.82 \pm 0.19$ \\

MetaCell2
& $0.10 \pm 0.13$
& $0.19 \pm 0.16$
& $0.63 \pm 0.59$
& $0.88 \pm 0.21$
& $0.05 \pm 0.10$
& $0.10 \pm 0.13$
& $0.44 \pm 0.44$
& $0.80 \pm 0.19$
& $0.09 \pm 0.07$
& $0.20 \pm 0.12$
& $0.63 \pm 0.54$
& $\underline{0.90 \pm 0.15}$ \\

SuperCell
& $0.14 \pm 0.12$
& $0.27 \pm 0.19$
& $0.38 \pm 0.43$
& $0.93 \pm 0.14$
& $0.07 \pm 0.10$
& $0.14 \pm 0.17$
& $0.21 \pm 0.28$
& $0.86 \pm 0.15$
& $0.10 \pm 0.09$
& $0.21 \pm 0.15$
& $0.73 \pm 0.50$
& $0.90 \pm 0.15$ \\

MetaQ
& $0.11 \pm 0.07$
& $0.24 \pm 0.13$
& $0.35 \pm 0.41$
& $0.96 \pm 0.11$
& $0.04 \pm 0.04$
& $0.08 \pm 0.08$
& $0.60 \pm 0.51$
& $0.81 \pm 0.16$
& $0.09 \pm 0.07$
& $0.21 \pm 0.14$
& $0.48 \pm 0.50$
& $\mathbf{0.91 \pm 0.13}$ \\
\bottomrule
\end{tabular}%
}
\setlength{\tabcolsep}{6pt}
\end{table*}

\subsection{Cross-Batch Aggregation Improves Coverage; Cohesion Preserves Homogeneity}

We examine whether density-corrected affinity combined with cross-batch aggregation improves representation of rare cell types.

As shown in Table~\ref{tab:rare_cells}, scProto's advantage is clearest against the batch-restricted metacell baselines (SEACells, SuperCell) on Pancreas and Lung, where a rare cell type often lacks enough cells within any single batch to form its own coherent metacell: on Pancreas, scProto leads both on homogeneity (0.48 vs.\ 0.42/0.14) and by a wide margin on F1 and coverage (0.54/0.77 vs.\ 0.37/0.43 for SEACells, 0.07/0.12 for SuperCell); on Lung the gap is smaller but in the same direction (homogeneity 0.48 vs.\ 0.47/0.45; F1 0.55 vs.\ 0.42/0.47; coverage 0.76 vs.\ 0.52/0.65). On Immune, where every batch already contains enough cells of each rare type for within-batch methods to reach full coverage, scProto is comparable rather than ahead: homogeneity (0.83), F1 (0.86), and coverage (1.00) all sit within 0.02 of SEACells and SuperCell's own scores. MetaCell2 and MetaQ follow a different pattern: both trail scProto clearly on Pancreas and Immune, but on Lung both score higher on homogeneity and F1 (MetaCell2: 0.56/0.64; MetaQ: 0.54/0.58; vs.\ scProto's 0.48/0.55); MetaCell2 also leads on Lung coverage (0.85 vs.\ scProto's 0.76), while MetaQ trails there (0.70).

Against the batch-correction baselines (scPoli-Stage-1+Leiden, Harmony, ComBat, BBKNN), scProto leads Harmony, ComBat, and BBKNN on every rare-cell metric on all three datasets --- consistent with these methods correcting batch effects at the level of a whole embedding or cluster, which can absorb a rare, batch-thin population into a denser, spuriously similar neighbor before any rare-cell signal is scored. scPoli-Stage-1+Leiden is the exception: scProto still leads it on Pancreas (homogeneity 0.48 vs.\ 0.45, F1 0.54 vs.\ 0.33, coverage 0.77 vs.\ 0.40) and on Immune (homogeneity 0.83 vs.\ 0.78, F1 0.86 vs.\ 0.82, coverage tied at 1.00), but trails it on Lung (homogeneity 0.48 vs.\ 0.57, F1 0.55 vs.\ 0.62, coverage 0.76 vs.\ 0.79) --- the one dataset where scProto's own pretrained encoder, clustered directly, out-recovers rare states relative to the full model.

\begin{table*}[t]
\centering
\caption{Rare-cell representation, pooled over train~+~held-out cells (cf.\ Table~\ref{tab:graph_batch}'s held-out-only protocol; metric definitions in \S\ref{sec:experiments}). scProto/scPoli (Stage-1) use the node-holdout model; other methods are fit on the complete dataset (context, not matched). Bold: best; underlined: second-best.}
\label{tab:rare_cells}
\tiny
\setlength{\tabcolsep}{1pt}
\resizebox{\textwidth}{!}{%
\begin{tabular}{lccc|ccc|ccc}
\toprule
 & \multicolumn{3}{c}{Pancreas} & \multicolumn{3}{c}{Immune} & \multicolumn{3}{c}{Lung} \\
\cmidrule(lr){2-4} \cmidrule(lr){5-7} \cmidrule(lr){8-10}
Method & Homog $\uparrow$ & F1 $\uparrow$ & Coverage $\uparrow$ & Homog $\uparrow$ & F1 $\uparrow$ & Coverage $\uparrow$ & Homog $\uparrow$ & F1 $\uparrow$ & Coverage $\uparrow$ \\
\midrule

scProto
& $\mathbf{0.48 \pm 0.10}$
& $\mathbf{0.54 \pm 0.11}$
& $\mathbf{0.77 \pm 0.20}$
& $0.83 \pm 0.08$
& $0.86 \pm 0.02$
& $\mathbf{1.00 \pm 0.00}$
& $0.48 \pm 0.19$
& $0.55 \pm 0.21$
& $0.76 \pm 0.28$ \\

scPoli (Stage-1) + Leiden
& $\underline{0.45 \pm 0.20}$
& $0.33 \pm 0.24$
& $0.40 \pm 0.24$
& $0.78 \pm 0.15$
& $0.82 \pm 0.16$
& $\underline{1.00 \pm 0.00}$
& $\mathbf{0.57 \pm 0.20}$
& $\underline{0.62 \pm 0.19}$
& $\underline{0.79 \pm 0.25}$ \\

\midrule

SEACells
& $0.42 \pm 0.20$
& $\underline{0.37 \pm 0.25}$
& $\underline{0.43 \pm 0.27}$
& $\mathbf{0.85 \pm 0.07}$
& $\underline{0.88 \pm 0.08}$
& $1.00 \pm 0.00$
& $0.47 \pm 0.24$
& $0.42 \pm 0.22$
& $0.52 \pm 0.23$ \\

BBKNN
& $0.16 \pm 0.06$
& $0.10 \pm 0.08$
& $0.17 \pm 0.13$
& $0.57 \pm 0.23$
& $0.60 \pm 0.29$
& $0.70 \pm 0.27$
& $0.36 \pm 0.11$
& $0.31 \pm 0.17$
& $0.45 \pm 0.24$ \\

Harmony
& $0.22 \pm 0.09$
& $0.22 \pm 0.17$
& $0.40 \pm 0.24$
& $0.64 \pm 0.13$
& $0.75 \pm 0.17$
& $0.93 \pm 0.13$
& $0.34 \pm 0.08$
& $0.28 \pm 0.14$
& $0.43 \pm 0.26$ \\

ComBat
& $0.23 \pm 0.10$
& $0.15 \pm 0.07$
& $0.33 \pm 0.27$
& $0.75 \pm 0.18$
& $0.80 \pm 0.19$
& $0.93 \pm 0.13$
& $0.40 \pm 0.21$
& $0.38 \pm 0.19$
& $0.53 \pm 0.26$ \\

MetaCell2
& $0.20 \pm 0.23$
& $0.20 \pm 0.26$
& $0.26 \pm 0.33$
& $0.44 \pm 0.28$
& $0.52 \pm 0.24$
& $0.78 \pm 0.19$
& $\underline{0.56 \pm 0.20}$
& $\mathbf{0.64 \pm 0.16}$
& $\mathbf{0.85 \pm 0.20}$ \\

SuperCell
& $0.14 \pm 0.06$
& $0.07 \pm 0.11$
& $0.12 \pm 0.18$
& $\underline{0.84 \pm 0.04}$
& $\mathbf{0.88 \pm 0.05}$
& $1.00 \pm 0.00$
& $0.45 \pm 0.16$
& $0.47 \pm 0.16$
& $0.65 \pm 0.22$ \\

MetaQ
& $0.20 \pm 0.12$
& $0.11 \pm 0.14$
& $0.25 \pm 0.33$
& $0.60 \pm 0.12$
& $0.73 \pm 0.13$
& $0.95 \pm 0.10$
& $0.54 \pm 0.22$
& $0.58 \pm 0.23$
& $0.70 \pm 0.24$ \\
\bottomrule
\end{tabular}%
}
\setlength{\tabcolsep}{6pt}
\end{table*}

\subsection{scProto captures microenvironmental structure while preserving cell-type identity}

We evaluate whether scProto's cell-type identity is grounded in gene expression itself rather than in whatever affinity graph it is given, unlike SEACells, which partitions whatever graph it is handed regardless of whether that structure is supported by gene expression. Using the BANKSY affinity's tunable spatial weight $\lambda$ (Appendix~\ref{app:spatial_affinity}), we move the graph continuously from purely transcriptomic ($\lambda{=}0$) to increasingly spatially dominated ($\lambda{=}0.2,0.4,0.8$) and train both methods on the identical graph at each $\lambda$. As shown in Fig.~\ref{fig:spatial_tradeoff}, SEACells' cell-type purity collapses as $\lambda$ rises ($0.82\to0.51$), and its niche-program recovery (Pearson $r$ between metacell- and single-cell-level logFC) degrades similarly ($0.85\to0.66$), since its kernel is built directly from the BANKSY-mixed embedding and simply groups whatever the graph says is close; scProto degrades far less on both (purity $0.87\to0.74$, recovery $0.81\to0.74$), because its identity comes from an encoder that reads the cell's own gene expression directly, $z_i = f_\theta(x_i, b_i)$, trained via reconstruction to stay faithful to it---so the graph only shapes how prototypes group already expression-grounded encodings, rather than determining identity itself. The same pattern holds at the gene level: of the ground-truth niche-DE genes each method's metacells recover (size-weighted Welch $t$-test, $p<0.05$, correct direction), scProto recovers \textbf{CXCL9} within Fibroblasts at the Desmoplastic-stroma niche in 3 of 4 graphs while SEACells recovers it in none---a literature-consistent, IFN$\gamma$-induced chemokine marking the tumor--stromal interface and recruiting CXCR3$^+$ T cells~\citep{chow2019cxcr3}, reported as one illustrative example rather than a systematic discovery claim.

\begin{figure*}[t]
\centering
\includegraphics[width=\textwidth]{figures/s28nsc_lambda_degradation.png}
\caption{
scProto vs.\ SEACells as the BANKSY spatial weight $\lambda$ increases (0 = no spatial signal, using the adaptive-RBF affinity on PCA embeddings only). Left: size-weighted cell-type purity. Right: niche-program recovery (Pearson $r$ between metacell-level and single-cell logFC, macro-averaged over cell-type/niche pairs). SEACells degrades sharply as $\lambda$ rises; scProto stays comparatively flat.
}
\label{fig:spatial_tradeoff}
\end{figure*}
